\documentclass[10pt,a4paper]{amsart}
\usepackage[margin=19mm]{geometry}
\usepackage{amsmath,amssymb,mathtools}
\usepackage[scr=rsfs,cal=euler]{mathalfa}
\usepackage{xfrac}
\usepackage[unicode,hidelinks,bookmarks=false]{hyperref}
\newcommand{\ie}{i.e.\ }
\newcommand{\cf}{cf.\ }
\newcommand{\resp}{respectively\ }
\newcommand{\Subsection}{\subsection}
\providecommand{\nobreakdash}{\nobreak\hbox{-}}
\setlength{\emergencystretch}{2em}
\multlinegap0pt
\usepackage{amssymb}
\usepackage{mathtools}
\usepackage[shortlabels]{enumitem}
\newtheorem{theorem}{Theorem}
\newcommand{\calE}{\mathcal E}
\newcommand{\calH}{\mathcal H}
\newcommand{\calL}{\mathcal L}
\newcommand{\calX}{\mathcal X}
\newcommand{\dxdt}{\,dx\,dt}
\newcommand{\harm}{\mathrm{harm}}
\newcommand{\loc}{\mathrm{loc}}
\newcommand{\nabh}{\nabla_h}
\title{Finite-Scale One-Component Regularity via Harmonic Pressure for the 3D Navier--Stokes Equations}
\author{Runlong Yu}
\date{Source: arXiv:2606.08352v1 (CC BY 4.0)}
\begin{document}
\maketitle

\noindent\emph{Source and licence.} Finite-Scale One-Component Regularity via Harmonic Pressure for the 3D Navier--Stokes Equations. Version arXiv:2606.08352v1 (CC BY 4.0), distributed by arXiv under the Creative Commons Attribution 4.0 International licence, http://creativecommons.org/licenses/by/4.0/, as recorded in the arXiv licence metadata for this exact version and shown on the abstract page https://arxiv.org/abs/2606.08352v1. Full licence notice: \texttt{licenses/arxiv-2606.08352-CC-BY-4.0.txt}.\par\smallskip

\noindent\textit{Editorial scope.} Counted unit: the article's theorem thm:qualitative-approx (source0.tex lines 720--734). The article's complete original proof of it is delivered with it (source0.tex lines 736--831, one \texttt{proof} environment of 96 lines). No mathematics is written, repaired or re-derived: the statement and the proof are the article's own lines, in the article's own notation, and no retained line is substituted or re-flowed.  Retained uncounted as the source-local context the proof needs: lines 96--100.  Nothing was omitted from any retained span, so the package carries no omission record.  No retained line contains a citation, so the wrapper carries no bibliography and no bibliography entry.  No retained line contains a figure environment, an image-inclusion command or drawing code, so no image is delivered and the package declares no asset dependency.  Editorial formatting only, in addition: an AMS \texttt{amsart} preamble that restates the article's own declarations and macros used by the retained lines, namely the environments \texttt{theorem} on separate counters, together with the standard packages those lines need.  The retained environments print in this document as Theorem 1 in order of appearance, whereas the article prints the same items with the numbering of its own full document.  The complete native source of the article as distributed in the arXiv e-print for this exact version is bundled unchanged as \texttt{sources/202325/source0.tex}.
\begin{equation}
        \partial_tu-\Delta u+(u\cdot\nabla)u+\nabla p=0,
        \qquad \nabla\cdot u=0.
        \label{eq:NS}
\end{equation}

\begin{theorem}[Qualitative finite-scale approximation]
\label{thm:qualitative-approx}
Let \(M\ge1\) and \(0<\theta<1/2\).  There exists a nondecreasing modulus
\[
        \omega_{M,\theta}:[0,\infty)\to[0,\infty),
        \qquad
        \lim_{s\downarrow0}\omega_{M,\theta}(s)=0,
\]
such that every suitable weak solution \((u,p)\) of \eqref{eq:NS} in \(Q_1\) satisfying \(\Phi(1)\le M\) obeys
\begin{equation}
        \calX^{\harm}_\theta(u,p;M)
        \le \omega_{M,\theta}(C_3(1)).
        \label{eq:qualitative-approx}
\end{equation}
\end{theorem}

\begin{proof}
Suppose the conclusion fails.  Then there exist \(\eta_0>0\) and suitable weak solutions \((u^{(n)},p^{(n)})\) in \(Q_1\) such that
\[
        \Phi_{u^{(n)}}(1)\le M,
        \qquad
        C^{(n)}_3(1)\to0,
\]
and
\[
        \calX^{\harm}_\theta(u^{(n)},p^{(n)};M)\ge\eta_0
        \qquad\text{for all }n.
\]
The uniform energy bound gives, after passing to a subsequence,
\[
        u^{(n)}\rightharpoonup v\quad\text{weakly in }L^2_tH^1_x(Q_\sigma),
        \qquad
        u^{(n)}\stackrel{*}{\rightharpoonup} v\quad\text{in }L^\infty_tL^2_x(Q_\sigma)
\]
for every \(\sigma<1\).  After fixing pressure representatives by subtracting functions of time, the equation gives a uniform bound for \(\partial_tu^{(n)}\) in
\[
        L^{3/2}((-\sigma^2,0);W^{-1,3/2}(B_\sigma)).
\]
Aubin--Lions gives strong convergence in \(L^2(Q_\sigma)\).  Interpolation with the uniform \(L^{10/3}\)-bound yields
\begin{equation}
        u^{(n)}\to v\quad\text{strongly in }L^3_{\loc}(Q_1).
        \label{eq:strong-L3}
\end{equation}
Since \(u^{(n)}_3\to0\) strongly in \(L^3(Q_1)\), the limit has the form \(v=(v_h,0)\), and \(\nabh\cdot v_h=0\).

Set
\[
        \widetilde p^{(n)}=p^{(n)}-(p^{(n)})_{B_1}(t).
\]
After extraction,
\[
        \widetilde p^{(n)}\rightharpoonup q
        \quad\text{weakly in }L^{3/2}_{\loc}(Q_1).
\]
Fix \(\rho\) with \(\theta<\rho<1\), and choose \(\chi\in C_c^\infty(B_\rho)\) equal to one on a neighborhood of \(B_\theta\).  In \(Q_\theta\), decompose
\[
        \widetilde p^{(n)}=P^{(n)}+H^{(n)},
        \qquad
        P^{(n)}=R_iR_j(\chi u_i^{(n)}u_j^{(n)}),
\]
where \(H^{(n)}(\cdot,t)\) is harmonic in \(B_\theta\) for a.e. time.  By \eqref{eq:strong-L3},
\[
        P^{(n)}\to P:=R_iR_j(\chi v_iv_j)
        \quad\text{strongly in }L^{3/2}(Q_\theta).
\]
Because \(\widetilde p^{(n)}\rightharpoonup q\) and \(P^{(n)}\to P\) strongly, \(H^{(n)}=\widetilde p^{(n)}-P^{(n)}\) converges weakly in \(L^{3/2}(Q_\theta)\), after extraction, to \(H=q-P\).  The limit is harmonic in space.  Set
\[
        h^{(n)}=H^{(n)}-H+(p^{(n)})_{B_1}(t).
\]
Then \(h^{(n)}\in\calH(Q_\theta)\), and
\[
        p^{(n)}-q-h^{(n)}=P^{(n)}-P\to0
        \quad\text{strongly in }L^{3/2}(Q_\theta).
\]
Together with \eqref{eq:strong-L3}, this gives
\[
        \calE^{\harm}_\theta((u^{(n)},p^{(n)}),(v,q);h^{(n)})\to0.
\]
It remains to identify the limit.  Passing to the limit in the horizontal momentum equation gives
\[
        \partial_t v_h-\Delta v_h+(v_h\cdot\nabh)v_h+\nabh q=0.
\]
The vertical equation gives, in distributions,
\[
        \partial_3p^{(n)}
        =-\partial_tu^{(n)}_3+\Delta u^{(n)}_3-\nabla\cdot(u^{(n)}u^{(n)}_3),
\]
and the right-hand side tends to zero in distributions using \(u^{(n)}_3\to0\) in \(L^3\), the uniform energy bounds, and \eqref{eq:strong-L3}.  Hence \(\partial_3q=0\).

The suitability of the limit follows from the standard stability of the local energy inequality under this convergence.  Indeed, for every fixed nonnegative test function \(\phi\in C_c^\infty(Q_\theta)\), the terms containing \(|u^{(n)}|^2\) and \(|u^{(n)}|^2u^{(n)}\) pass to the limit by the strong \(L^3\)-convergence.  The pressure term also passes to the limit, since \(\widetilde p^{(n)}\rightharpoonup q\) in \(L^{3/2}\) and \(u^{(n)}\to v\) strongly in \(L^3\), so
\[
        \int \widetilde p^{(n)}u^{(n)}\cdot\nabla\phi\,\dxdt
        \longrightarrow
        \int qv\cdot\nabla\phi\,\dxdt .
\]
The subtracted time-dependent pressure averages do not contribute because \(u^{(n)}\) is divergence-free.  Lower semicontinuity gives the dissipative term.  Thus \((v,q)\in\calL_M(Q_\theta)\), after increasing \(K_0(M,\theta)\) if necessary.  This is incompatible with the lower bound on \(\calX^{\harm}_\theta\).

Define
\[
\omega_{M,\theta}(\delta)
=\sup\{\calX^{\harm}_\theta(u,p;M): (u,p)\text{ suitable in }Q_1,
\ \Phi(1)\le M,
\ C_3(1)\le\delta\}.
\]
The preceding contradiction argument proves \(\omega_{M,\theta}(\delta)\to0\) as \(\delta\downarrow0\).  The supremum is finite by a direct fixed-scale comparison: choose \((v,q)=(0,0)\in\calL_M(Q_\theta)\) and choose the harmonic corrector \(h=(p)_{B_\theta}(t)\), which is constant in space.  Then
\[
        \calX^{\harm}_\theta(u,p;M)
        \le C_\theta\{C(\theta)+D(\theta)\}
        \le C_\theta M,
\]
using fixed-scale nesting.  Replacing \(\omega_{M,\theta}\) by its monotone envelope gives a nondecreasing modulus.
\end{proof}

\end{document}
